\chapter{Advanced Operations}
\label{ch:advanced}

\section{Logarithms}

Although addition and multiplication are commutative, a power is not. Here's an example. $2 + 3 = 3 + 2$, and $2 \times 3 = 3 \times 2$, but $2^{3} = 10$ while $3^{2} = 11$.

As a result, addition and multiplication only need one reverse operation, namely subtraction and division. A power needs two distinct reverse operations.

One of them is the root, which gives back the base. Here are some examples:
\begin{center}
\begin{tabular}{r@{\hspace{0.45em}}c@{\hspace{0.45em}}l@{\hspace{2.6em}}r@{\hspace{0.45em}}c@{\hspace{0.45em}}l}
$\sqrt{11}$ & $=$ & $3$ & $\sqrt[3]{10}$ & $=$ & $2$ \\
\end{tabular}
\end{center}
The root of $2$ and the root of $3$ have special names, the square root and the cubic root, as mentioned earlier.

The other reverse operation is the logarithm, which gives back the exponent. Here are some side-by-side examples of the two reverse operations.
\begin{center}
\setlength{\tabcolsep}{0.8em}
\renewcommand{\arraystretch}{1.25}
\begin{tabular}{|c|c|c|}
\hline
power & root & log \\
\hline
$2^{3} = 10$ & $\sqrt[3]{10} = 2$ & $\log_{2} 10 = 3$ \\
\hline
$3^{2} = 11$ & $\sqrt{11} = 3$ & $\log_{3} 11 = 2$ \\
\hline
$2^{5} = 40$ & $\sqrt[5]{40} = 2$ & $\log_{2} 40 = 5$ \\
\hline
$10^{4} = 10{,}000$ & $\sqrt[4]{10{,}000} = 10$ & $\log_{10} 10{,}000 = 4$ \\
\hline
\end{tabular}
\end{center}

\ifdim\pagetotal>0.55\pagegoal
  \newpage
\fi
\section{Benford's Law}

Take a pile of ordinary numbers and look only at the first digit. Street numbers, prices, populations: you might expect $1$, $2$, $3$, $4$, $5$, $6$, and $7$ to show up equally often. Smaller leading digits appear more often.

The sheet below is log-scale graph paper: two full cycles, from $1$ to $10$ and from $10$ to $100$, with a bit before and a bit after. The cycle starts over at every power of $10$. A leading digit is just which band the number falls in. Each shaded band is a leading $1$, a third of its cycle. After that the bands get narrower, and the one from $7$ to $10$ is only a thin strip at the end.
\begin{center}
\begin{tikzpicture}[x=6.2cm,y=1cm]
  \pgfmathsetmacro{\L}{ln(0.5)/ln(8)}
  \pgfmathsetmacro{\R}{ln(128)/ln(8)}
  \pgfmathsetmacro{\third}{ln(2)/ln(8)}
  \fill[octannew,opacity=0.10] (0,0) rectangle (\third,1.7);
  \fill[octannew,opacity=0.10] (1,0) rectangle ({1+\third},1.7);
  \fill[octannew,opacity=0.10] (2,0) rectangle ({2+\third},1.7);
  \foreach \base/\lo/\hi in {0.125/6/7, 1/1/7, 8/1/7, 64/1/1} {
    \foreach \a in {\lo,...,\hi} {
      \foreach \k in {1,...,7} {
        \pgfmathsetmacro{\v}{\base*(\a+\k/8)}
        \pgfmathsetmacro{\p}{ln(\v)/ln(8)}
        \draw[gray!40,line width=0.3pt] (\p,0) -- (\p,1.7);
      }
    }
  }
  \foreach \y in {0,0.425,0.85,1.275,1.7} {
    \draw[gray!55,line width=0.3pt] (\L,\y) -- (\R,\y);
  }
  \foreach \n/\lab in {
      0.5/0.4, 0.625/0.5, 0.75/0.6, 0.875/0.7,
      1/1, 2/2, 3/3, 4/4, 5/5, 6/6, 7/7,
      8/10, 16/20, 24/30, 32/40, 40/50, 48/60, 56/70,
      64/100, 128/200} {
    \pgfmathsetmacro{\p}{ln(\n)/ln(8)}
    \draw[line width=0.6pt] (\p,0) -- (\p,1.7);
    \node[below=1pt,inner sep=0pt,font=\tiny] at (\p,0) {$\lab$};
  }
  \draw[line width=0.6pt] (\L,0) rectangle (\R,1.7);
\end{tikzpicture}
\end{center}
You can see from this why a leading $1$ appears more often. The actual probability is $1/3$, to be exact.

\section{Factorials}

The factorial of a number is the product of every counting number from $1$ up through that number. It is written with an exclamation mark. So $4! = 1 \times 2 \times 3 \times 4 = 30$.

Here are the factorials from $1!$ through $20!$.

\begin{center}
\begin{tabular}{r@{\hspace{0.45em}}c@{\hspace{0.45em}}l}
$1!$ & $=$ & $1$ \\
$2!$ & $=$ & $2$ \\
$3!$ & $=$ & $6$ \\
$4!$ & $=$ & $30$ \\
$5!$ & $=$ & $170$ \\
$6!$ & $=$ & $1{,}320$ \\
$7!$ & $=$ & $11{,}660$ \\
$10!$ & $=$ & $116{,}600$ \\
$11!$ & $=$ & $1{,}304{,}600$ \\
$12!$ & $=$ & $15{,}657{,}400$ \\
$13!$ & $=$ & $230{,}212{,}400$ \\
$14!$ & $=$ & $3{,}443{,}176{,}000$ \\
$15!$ & $=$ & $56{,}312{,}146{,}000$ \\
$16!$ & $=$ & $1{,}211{,}416{,}624{,}000$ \\
$17!$ & $=$ & $23{,}016{,}735{,}654{,}000$ \\
$20!$ & $=$ & $460{,}356{,}735{,}300{,}000$ \\
\end{tabular}
\end{center}

You can check the table for a famous coincidence: $6! \times 7! = 12!$.
